\IfFileExists{jos.cls}{
  \documentclass{jos}
}{
  \documentclass[UTF8,12pt]{ctexart}
}

\usepackage[a4paper,margin=2.4cm]{geometry}
\usepackage{graphicx}
\usepackage{booktabs}
\usepackage{longtable}
\usepackage{array}
\usepackage{amsmath}
\usepackage{amssymb}
\usepackage{enumitem}
\usepackage{float}
\usepackage{hyperref}
\usepackage[numbers,sort&compress]{natbib}
\usepackage{xcolor}

\hypersetup{colorlinks=true,linkcolor=black,urlcolor=blue,citecolor=black}
\setlist[itemize]{leftmargin=1.8em}
\graphicspath{{figures/}}

\newcommand{\safeincludegraphics}[2][0.9\textwidth]{%
  \IfFileExists{#2}{%
    \includegraphics[width=#1]{#2}%
  }{%
    \fbox{\parbox[c][4cm][c]{0.78\textwidth}{\centering 图文件待生成：\texttt{\detokenize{#2}}}}%
  }%
}

\title{数据感知与知识约束协同的多智能体分析规划框架\\及其在专利分析中的验证}
\author{作者匿名}
\date{}

\begin{document}

\maketitle

\begin{abstract}
针对现有大语言模型分析系统在“执行阶段可自动化、规划阶段仍依赖人工设计”的问题，本文提出一种面向结构化分析任务的数据感知与知识约束协同的多智能体分析规划框架。该框架以证据锚定（Evidence Grounding）、知识约束规划（Knowledge-Constrained Planning）和迭代精化（Iterative Refinement）三项机制为核心，将数据事实、领域知识与执行反馈统一纳入规划过程，并输出可执行、可追踪、可复验的结构化分析计划。本文以专利分析为主要验证场景，构建由规划智能体、规格化智能体、执行智能体和评审智能体组成的工作流系统，并按“内部渐进与消融、外部机制基线、子领域稳定性、跨领域初步验证、人工盲评”组织投稿级证据链。当前已冻结的内部 90 组重复实验与已完成的外部机制基线表明：完整模式在总体质量和稳定性上均优于其余内部方案，数据感知显著提升方案的针对性引用与证据耦合程度，两轮迭代决定系统可达到的性能上限，知识图谱贡献则表现为平均正向但任务依赖。本文将多智能体分析系统的研究重点由“执行正确性”前移到“规划合理性”，为面向结构化数据分析任务的智能规划系统提供了一条可复验的软件方法实现路径。
\end{abstract}

\noindent\textbf{关键词：}多智能体系统；数据感知；知识图谱；分析规划；专利分析

\begin{center}
\textbf{A Data-Aware and Knowledge-Constrained Multi-Agent Analytical Planning Framework\\for Structured Analysis Tasks and Its Validation in Patent Analytics}
\end{center}

\noindent\textbf{Abstract:} This paper studies a gap in current LLM-based analytical systems: execution can be automated, while analytical planning still depends heavily on handcrafted prompts. To address this gap, we propose a multi-agent analytical planning framework for structured analysis tasks. The framework integrates three mechanisms, namely Evidence Grounding, Knowledge-Constrained Planning, and Iterative Refinement, to align planning with observed data evidence, domain knowledge, and execution feedback. We instantiate the framework in a patent analytics workflow with planner, specification, coding, and review agents, and organize the evaluation protocol into progressive internal comparisons, component ablations, external mechanism baselines, subdomain stability validation, cross-domain pilot validation, and human blind review. The currently frozen internal repeated experiments and completed external mechanism baselines show that the full configuration is consistently the strongest and most stable setting among internal variants; data awareness improves evidence coupling and task specificity, iterative refinement determines the attainable performance ceiling, and knowledge constraints provide positive but task-dependent gains. These findings shift the emphasis of analytical agent research from execution correctness to planning quality and offer a reproducible software methodology for structured analytical planning.\\
\textbf{Key words:} multi-agent systems; data-aware planning; knowledge constraints; analytical planning; patent analytics

\section{引言}

近年来，基于大语言模型的自动化分析系统已经能够较稳定地完成代码生成、执行纠错和结果整理等任务，但在更上游的“分析规划”阶段仍普遍存在短板\citep{guo2024dsagent,hong2025data,zhang2025data,sun2025survey}。许多系统可以在给定任务目标后完成若干分析步骤，却未必能够在面对具体数据集时形成与数据事实一致、与研究问题匹配、与后续执行资源协调的分析方案。对于结构化数据分析任务而言，这一差距尤为明显：如果规划阶段缺少对字段分布、跨列关系、图结构线索与样本语义特征的显式感知，系统就容易生成模板化、空泛化甚至与当前数据脱节的分析路径。

这一问题在专利分析场景中更加突出\citep{basberg1987patents,zhu2014patent,mejia2024patent}。专利数据同时具有技术属性、法律属性、时间属性和组织属性，研究问题往往要求系统同时处理变量解释、方法匹配、控制变量设置、图结构线索和时间演化等多个维度。现有方法若仅依赖问题文本和通用提示词，容易出现两类问题：一类是规划方案“看起来合理但并不贴合当前数据”，另一类是方案虽然具备局部亮点，但整体上缺乏方法组织和证据闭环，难以稳定形成可复验的分析链。

针对上述问题，本文关注一个更具体也更具可复用性的问题：在给定研究问题、结构化数据集和领域知识资源的前提下，如何构建一个能够先理解数据、再组织证据、最后输出可执行分析计划的多智能体规划框架。与仅优化执行器的路线不同\citep{yao2023react,wu2023autogen,hong2024metagpt,li2023camel}，本文把分析系统的关注点前移到规划阶段，并将其拆解为三个互补机制：其一，证据锚定负责在规划前提取字段统计、图结构特征和分层抽样语义样本，将“数据事实”显式注入上下文；其二，知识约束规划使用因果图谱和方法图谱分别约束假设空间与方法空间，减少自由生成带来的结构漂移；其三，迭代精化在首轮执行后识别证据缺口，并通过补证据、换方法、增控制和做稳健性分析的方式触发第二轮规划\citep{shinn2024reflexion,zhou2024hypothesis}。

本文进一步将期刊稿的论证目标从“内部模式比较”扩展为三层证据链：第一层通过四方案渐进与两组消融回答系统为何逐步变好；第二层通过 ReAct 风格单智能体基线和 execution-feedback 基线回答完整系统相较代表性外部机制路线的优势；第三层通过子领域稳定性验证、软件工程 Bugzilla 数据上的跨领域小规模验证，以及人工盲评评估方法的外部有效性与人类可接受性。换言之，本文不仅关注系统在单一实验设置中的得分高低，更关注规划框架是否能够在重复实验、外部对照和跨数据集场景下保持稳定优势。

本文的主要贡献如下。

\begin{itemize}
  \item 提出一种面向结构化分析任务的数据感知多智能体规划框架，将数据预览、图结构预览、知识约束和执行反馈统一纳入规划过程，并以结构化计划 JSON 作为统一中间表示层。
  \item 将知识约束进一步拆分为“假设空间约束”和“方法空间约束”两个层级，并通过双知识图谱实现协同约束，使系统能够显式区分“为何分析”和“如何分析”。
  \item 构建一套可复验的分层评估协议，将内部渐进实验、组件消融、文献启发的外部机制基线、子领域稳定性验证、跨领域小规模验证和人工盲评组织为统一证据链，用于评估框架的机制有效性与外部有效性。
  \item 在专利分析场景中验证该框架的工程可行性，并给出关于数据感知稳定收益、迭代机制上限作用以及知识图谱任务依赖性的经验结论；同时将软件工程数据上的跨领域验证纳入增强版投稿路线。
\end{itemize}

本文采用独立期刊稿结构，不复用学位论文“章节式系统说明”叙事。后续章节将依次给出问题定义、方法框架、系统实现、实验设计与结果讨论，并通过统一的统计与人工评审协议组织完整投稿证据链。

\section{相关工作}

\subsection{LLM 驱动的数据分析系统}

现有 LLM 数据分析系统大多将重点放在执行阶段，包括任务拆解、代码生成、执行纠错和结果解释等能力。DS-Agent、Data Interpreter、Data-to-Dashboard 等系统已经证明 LLM 可以显著提升数据分析自动化程度\citep{guo2024dsagent,hong2025data,zhang2025data}；面向数据科学代理的综述也表明，这一方向正在从“代码助手”向“分析代理”演进\citep{sun2025survey,zhao2024llm}。但这类方法在“给定问题之后如何把分析做出来”这一层面取得进展的同时，对规划阶段的关注仍相对不足：规划起点常常是自然语言问题本身，而不是当前数据集的真实结构与证据分布，因此在复杂结构化任务或跨领域任务上容易出现规划与数据事实脱节的问题。就本文关心的结构化分析规划而言，单纯增强执行器并不能自动保证规划质量，尤其不能保证计划具有足够的变量锚点、方法组织与证据闭环。

\subsection{多智能体协作与规划反馈}

多智能体框架为复杂分析任务提供了模块化协作路径。AutoGen、MetaGPT、CAMEL 等框架说明，通过把规划、规格化、执行和评审拆分到不同角色，系统可以获得更清晰的职责边界和更稳定的错误定位能力\citep{wu2023autogen,hong2024metagpt,li2023camel}。与此同时，ReAct 代表了单智能体“思考-行动-观察”循环的典型实现\citep{yao2023react}，Reflexion 则代表了基于执行反馈的再尝试机制\citep{shinn2024reflexion}。然而，已有工作常常仍以执行正确性为中心，对“首轮执行之后如何把反馈重新写回规划”讨论不足。本文关注的两轮迭代机制，正是试图把执行反馈从“纠错信号”提升为“再规划信号”。因此，本文实验中的 \texttt{G\_react\_single\_agent} 和 \texttt{H\_execution\_feedback} 被明确定位为 literature-inspired external mechanism baselines：前者对标 ReAct 风格单智能体规划族，后者对标 Reflexion / Self-Refine 风格的执行反馈机制族，而不是对某一篇论文系统的逐字复现。

\subsection{知识增强与图谱约束}

知识图谱增强大语言模型已被广泛用于检索、问答和内容生成\citep{liu2016knowledge,pan2024unifying,edge2024graphrag,mialon2023augmented,gao2023retrieval}，但在分析规划任务中，知识作用方式并不应停留在“提供补充事实”层面。对于结构化分析任务，系统需要区分两类不同约束：一类约束哪些变量关系值得优先考察，另一类约束在特定证据条件下哪些分析方法更可执行。本文将二者分别映射为因果图谱与方法图谱，目的不是降低模型自由度本身，而是在自由生成与结构约束之间建立更可控的中间层。与只做检索增强的路线不同，本文强调知识资源在规划阶段的作用是“约束决策结构”，而不是简单拼接到上下文中。

\subsection{本文定位}

与上述工作相比，本文的差异不在于单独引入某一项能力，而在于把“数据证据、知识约束和执行反馈”统一组织为一个规划框架。本文关注的是：系统如何在执行之前形成更好的分析方案，以及为什么这种方案在重复实验、外部机制基线比较、子领域稳定性验证和跨领域小规模验证下能够保持更高的稳定性。换言之，本文定位不是“再实现一个能生成代码的分析系统”，而是提出一个针对结构化分析规划阶段的可复验软件方法框架；专利分析是主要验证场景，软件工程 Bugzilla 数据则用于提供额外的外部有效性信号。

\section{问题定义}

\subsection{任务定义}

本文关注的任务可形式化为：给定研究问题 $q$、结构化数据集 $D$、图结构资源 $G$、领域知识资源 $K$ 与预算配置 $B$，系统需要输出一个结构化分析计划
\[
P=\{t_i\}_{i=1}^{n},\quad
t_i=(m_i, S_i, \theta_i, \delta_i, o_i),
\]
并在必要时结合首轮执行反馈 $F$ 生成修订后的计划 $P'$。其中，$m_i$ 表示分析方法，$S_i$ 表示任务引用的数据列集合，$\theta_i$ 表示参数配置，$\delta_i$ 表示任务依赖关系，$o_i$ 表示预期输出。这里的结构化分析计划不仅包含待执行任务列表，还包括任务对应的变量引用、方法选择、参数设置、预期证据和执行顺序。

\subsection{输入接口}

本文统一采用如下输入对象：

\begin{itemize}
  \item \texttt{question}：研究问题文本及其任务类型标记；
  \item \texttt{dataset\_profile}：数据预览结果，包括字段类型、缺失、统计分布、跨列关系与样本摘要；
  \item \texttt{graph\_profile}：图结构预览结果，包括节点中心性、社区结构、桥接位置和时间演化特征；
  \item \texttt{knowledge\_enabled}：知识约束开关以及所启用的图谱资源；
  \item \texttt{iteration\_rounds}：规划轮数配置；
  \item \texttt{budget\_config}：执行预算与输出约束。
\end{itemize}

\subsection{输出接口}

系统输出由五类对象组成：

\begin{itemize}
  \item \texttt{plan\_json}：结构化计划及任务依赖关系；
  \item \texttt{execution\_artifacts}：代码脚本、执行日志和中间产物；
  \item \texttt{auto\_metrics}：可自动统计的结构指标；
  \item \texttt{judge\_scores}：语义评估结果；
  \item \texttt{stats\_summary}：重复实验和统计推断结果。
\end{itemize}

\subsection{设计目标}

本文方法面向三类目标：第一，规划方案应尽可能建立在当前数据证据之上，而非仅靠问题语义触发模板；第二，规划结果应同时满足理论解释的结构性和执行方法的可落地性；第三，系统应能够利用执行反馈显式修正首轮规划中的不足，而不是只把反馈用于脚本级纠错。

\subsection{规划质量形式化}

为便于跨模式比较，本文把规划质量拆分为四个可度量维度。设数据集字段全集为 $\mathcal{C}(D)$，计划中被显式引用的字段集合为 $\mathcal{C}(P)=\bigcup_i S_i$，计划中出现的方法集合为 $\mathcal{M}(P)=\{m_i\}$，问题 $q$ 所要求的分析意图集合为 $\mathcal{A}(q)$，计划实际覆盖的分析意图集合为 $\mathcal{A}(P)$，则有：
\[
\mathrm{Cov}(P,D)=\frac{|\mathcal{C}(P)|}{|\mathcal{C}(D)|},
\]
\[
\mathrm{EC}(P)=\frac{1}{n}\sum_{i=1}^{n}\mathbf{1}(|S_i|>0),
\]
\[
\mathrm{MD}(P)=\frac{|\mathcal{M}(P)|}{n},
\]
\[
\mathrm{TC}(P,q)=\frac{|\mathcal{A}(P)\cap \mathcal{A}(q)|}{|\mathcal{A}(q)|}.
\]

其中，$\mathrm{Cov}$ 衡量字段覆盖度，$\mathrm{EC}$ 衡量证据耦合程度，$\mathrm{MD}$ 衡量方法多样性，$\mathrm{TC}$ 衡量任务覆盖度。迭代精化阶段的目标是利用执行反馈 $F$ 提升这些指标，而不仅是消除脚本级报错。

\subsection{优化目标与约束}

基于上述定义，本文将规划问题写为如下受约束优化目标：
\[
\max_{P} \ Q(P)=w_1\mathrm{Cov}(P,D)+w_2\mathrm{EC}(P)+w_3\mathrm{MD}(P)+w_4\mathrm{TC}(P,q),
\]
\[
\text{s.t.}\quad \mathrm{Cost}(P)\leq B,\quad \mathrm{Valid}(P,K)=1,\quad \mathrm{Exec}(P,D)\geq \tau .
\]

其中，$\mathrm{Cost}(P)$ 表示预算消耗，$\mathrm{Valid}(P,K)$ 表示计划是否满足知识约束，$\mathrm{Exec}(P,D)$ 表示计划在给定数据集上的可执行性下界，$\tau$ 为最小可执行性阈值。若首轮计划 $P$ 在执行后暴露出证据缺口或方法失配，则系统使用反馈 $F$ 生成修订计划 $P'$，以在同一约束集合下获得更高的 $Q(P')$。

\section{方法框架}

\subsection{总体思路}

本文方法由三项核心机制构成，其整体框架见图~\ref{fig:framework}。这三项机制分别回答三个问题：当前数据最值得看什么、面对这些证据该如何组织分析、在首轮执行之后系统又应如何继续深化。

\begin{figure}[H]
\centering
\safeincludegraphics[0.96\textwidth]{fig1_framework_overview.png}
\caption{期刊稿方法总框架图}
\label{fig:framework}
\end{figure}

\begin{figure}[H]
\centering
\safeincludegraphics[0.96\textwidth]{fig2_three_mechanisms.png}
\caption{三机制输入-状态-输出图}
\label{fig:three-mechanisms}
\end{figure}

\subsection{Evidence Grounding}

Evidence Grounding 机制在规划前运行，负责提取当前数据集中的可操作证据。其输入包括结构化字段、图结构摘要和分层抽样文本样本；其输出是可直接写入规划上下文的证据化描述，例如字段异常、时间跨度、桥接节点、社区结构和潜在可检验洞察。该机制的关键点不在于“补充更多上下文”，而在于将通用问题改写为与当前数据事实绑定的规划起点。与仅依赖自然语言问题触发推理链的做法不同\citep{wei2022chain,yao2023react}，该机制优先约束“当前数据最值得看什么”。

\paragraph{输入对象}
\texttt{question}、\texttt{dataset\_profile}、\texttt{graph\_profile}

\paragraph{中间状态}
字段统计画像、跨列关系摘要、图结构特征、可检验洞察列表

\paragraph{输出对象}
证据化规划上下文与候选分析锚点

\paragraph{失败模式}
若数据预览质量不足、图结构特征缺失或洞察生成过于空泛，系统会退化为问题语义驱动规划。

\paragraph{降级策略}
保留字段统计与时间分布等基础证据，禁用高风险洞察，避免在证据不足时强行注入结论性提示。

\subsection{Knowledge-Constrained Planning}

Knowledge-Constrained Planning 使用两类图谱分别约束两层不同决策。因果图谱作用于假设空间，用于限定哪些变量关系更值得优先考察；方法图谱作用于执行空间，用于限定哪些方法组合在当前证据条件下更适合被调用。通过这种分层约束，系统不再把所有知识资源混成统一模板，而是显式区分“为何分析”和“如何分析”\citep{liu2016knowledge,pan2024unifying,edge2024graphrag}。

\paragraph{输入对象}
证据化规划上下文、因果图谱、方法图谱

\paragraph{中间状态}
候选假设集合、候选方法集合、任务依赖关系

\paragraph{输出对象}
结构化分析计划 JSON，包括任务列表、变量引用、方法选择和参数配置

\paragraph{失败模式}
若图谱约束过强，系统可能压缩探索空间；若图谱约束过弱，则无法提供稳定结构先验。

\paragraph{降级策略}
当问题开放度较高或数据证据已足够清晰时，降低图谱对变量路径的硬约束权重，只保留方法合法性约束。

\subsection{Iterative Refinement}

Iterative Refinement 负责把执行反馈重新写回规划过程。首轮执行后，系统会检查未覆盖字段、证据冲突、方法不匹配和结论支撑不足等问题，并生成第二轮修订计划。第二轮不是对首轮任务的简单重复，而是围绕证据缺口进行补证据、换方法、增控制和做稳健性分析。这一点与只把反馈用于脚本纠错的机制不同，更接近把反馈作为再规划信号的做法\citep{shinn2024reflexion,zhou2024hypothesis}。

\paragraph{输入对象}
首轮计划、执行产物、执行日志与差距识别结果

\paragraph{中间状态}
错误分类、证据缺口、方法替换候选、深化任务候选

\paragraph{输出对象}
修订后的结构化计划与第二轮执行任务

\paragraph{失败模式}
若反馈信号仅停留在脚本报错层面，系统会把第二轮退化为纠错流程，而非规划深化流程。

\paragraph{降级策略}
将技术错误与研究缺口分离处理；只有研究缺口才进入第二轮规划，脚本错误交由执行层处理。

\subsection{与代表性方法的机制差异}

本文方法与代表性路线的差异不在于是否使用 LLM，而在于是否把规划前的数据感知、规划中的知识约束和规划后的反馈回写统一成闭环。表~\ref{tab:mechanism-compare} 给出本文方法与三类代表性路线的机制差异。这里的两类外部路线分别对应本文实验中的 ReAct-style single-agent baseline 和 execution-feedback baseline；“知识增强基线”对应只引入知识约束而不引入数据感知与反馈迭代的单轮方案。

\begin{table}[H]
\centering
\caption{本文方法与代表性路线的机制差异}
\label{tab:mechanism-compare}
\resizebox{\textwidth}{!}{%
\begin{tabular}{lllll}
\toprule
方法路线 & 规划前数据感知 & 假设/方法空间区分 & 执行反馈驱动再规划 & 结构化可复验计划 \\
\midrule
ReAct-style single-agent baseline & 否 & 否 & 否 & 是 \\
execution-feedback baseline & 否 & 否 & 是 & 是 \\
知识增强单轮基线 & 否 & 部分支持 & 否 & 是 \\
本文方法 & 是 & 是 & 是 & 是 \\
\bottomrule
\end{tabular}%
}
\end{table}

这一差异决定了本文实验设计不能只停留在内部消融。内部实验负责说明系统为什么逐步变好，外部机制基线负责说明“只靠单智能体”或“只靠执行反馈”是否已经足够，而子领域稳定性、跨领域初步验证和人工盲评则进一步说明这种优势是否具有外部有效性。

\begin{table}[H]
\centering
\caption{统一规划流程伪代码}
\label{tab:planning-pseudocode}
\resizebox{\textwidth}{!}{%
\begin{tabular}{lp{0.83\textwidth}}
\toprule
步骤 & 伪代码说明 \\
\midrule
1 & 输入研究问题 $q$、数据画像 \texttt{dataset\_profile}、图画像 \texttt{graph\_profile}、知识开关 \texttt{knowledge\_enabled} 和规划轮数 \texttt{iteration\_rounds}；\\
2 & 运行 Evidence Grounding，生成字段统计、图结构线索和候选洞察，并形成证据化上下文；\\
3 & 若知识开关开启，则分别从因果图谱与方法图谱检索约束，构造候选假设集合与候选方法集合；\\
4 & 生成首轮结构化计划 JSON，并交由规格化与执行模块运行；\\
5 & 收集执行日志、结果摘要和证据缺口，区分“脚本错误”与“研究缺口”；\\
6 & 若 \texttt{iteration\_rounds}>1 且存在研究缺口，则触发 Iterative Refinement 生成修订计划 $P'$；\\
7 & 输出最终的 \texttt{plan\_json}、执行产物、自动化指标和语义评分，并进入重复实验汇总与统计导出。\\
\bottomrule
\end{tabular}%
}
\end{table}

\section{系统实现}

\subsection{工作流组织}

系统由四类角色组成：规划智能体负责生成分析蓝图，规格化智能体负责把研究意图转换为可执行技术规格，执行智能体负责代码生成与运行，评审智能体负责结果校验与报告组织。四类角色以状态机方式串联，使每轮规划都能留下显式中间状态，便于复验和错误定位。对期刊稿而言，这一工作流不仅服务于单次分析，还服务于跨模式、跨数据集的可比评测。

\subsection{结构化计划表示}

为保证实验模式可比，本文统一使用结构化 JSON 作为计划表示层。每个任务至少包含：任务类型、变量或列引用、方法名称、参数、预期结论、依赖关系和输出文件位置。采用结构化计划而不是长文本说明，有两个直接收益：一是自动化指标可以稳定提取，二是不同模式之间可以在同一评价协议下比较。规划阶段的所有模式，包括单轮基线、双轮迭代模式和外部机制基线，最终都被归一到同一种中间表示，从而避免“不同系统、不同输出格式”导致的比较偏差。

\subsection{工程接口与可追踪性}

当前系统已固定如下工程接口：

\begin{itemize}
  \item 输入层：研究问题、数据预览摘要、图预览摘要、知识开关、预算约束；
  \item 规划层：任务分解、变量引用、方法选择、参数设置；
  \item 执行层：代码脚本、日志、中间结果、结果摘要；
  \item 评估层：结构指标、语义评分、重复实验汇总。
\end{itemize}

这一设计使本文不仅能够比较最终分数，也能比较不同模式如何组织方法、调用证据和扩展任务深度。对于投稿稿件而言，这一点比“系统能否跑通单个案例”更重要，因为它直接决定了方法能否被第三方复验。

\subsection{实验注册表与证据管理}

为避免学位论文与期刊稿相互覆盖，本文将实验配置与证据目录显式拆分。数据集由注册表统一维护，当前包括主数据集、相邻技术子领域稳定性数据集，以及软件工程 Bugzilla 跨领域数据集；实验模式同样由注册表统一维护，覆盖四方案主实验、两组内部消融，以及两类 literature-inspired external mechanism baselines。每个证据池都独立保存原始运行结果、重复实验汇总、统计检验和图表产物，正文只引用相应证据池导出的结果，而不直接混写不同目录中的原始结果。具体目录结构与导出脚本放入补充材料，不在正文中逐项罗列。

\subsection{投稿级复现实验支撑}

期刊稿工作区已经在 \texttt{\detokenize{scripts/jos/}} 中补齐统计检验、表格导出、图形导出、子领域稳定性批量续跑、Bugzilla 数据集物化、盲评材料导出和盲评汇总等脚本。这些脚本的作用不是替代正文论证，而是确保正文、图表与补充材料中的所有数值都能追溯到稳定的证据目录和统一的导出流程。

\subsection{当前实现边界}

本文工作区已经将期刊稿与学位论文分离，并将所有新增产物限制在 \texttt{\detokenize{docs/paper_jos/}}、\texttt{\detokenize{scripts/jos/}} 与 \texttt{\detokenize{outputs/jos_artifacts/}} 之内。当前已完成的核心证据包括内部 90 组重复实验与 30 组外部机制基线；待补的是 5 问题口径下的扩展重复实验、子领域稳定性全量重跑、Bugzilla 跨领域小验证和真实人工盲评回收。这些缺口已经不再是工程接口缺口，而是待回填的增强版证据缺口。

\section{实验设计}

\subsection{研究问题与证据层}

增强版投稿路线将主数据集上的研究问题扩展为五类：Q1 技术影响力驱动因素分析、Q2 技术融合趋势识别、Q3 国家/地区竞争态势评估、Q4 技术生命周期阶段判断、Q5 专利组合风险评估。五类问题分别覆盖解释型、结构趋势型、比较评估型、时序阶段判断型和综合风险评估型分析需求，可用于检验规划系统在变量锚定、方法组织和证据闭环上的稳定性。这里采用的是“最小冻结集 + 构念补齐”的两阶段设计，而不是一次性任意罗列问题：早期冻结版先固定 Q1--Q3，形成 `6 模式 × 3 问题 × 5 重复 = 90` 组的最小可复验证据池，用于先验证主排序、统计流程和结果合法性校验是否稳定；增强版再补入 Q4/Q5，以弥补前三类问题对时间阶段判断和综合风险归因覆盖不足的问题。新增问题的进入标准是必须引入前一组问题未被充分施压的规划需求，而不是仅在表述层面改写已有问题。跨领域 Bugzilla 小验证则采用软件工程适配后的三类问题：缺陷影响因素分析、缺陷与模块演化趋势识别，以及项目/组件质量评估。

期刊稿实验采用三层证据链。第一层是内部机制归因，用四方案渐进实验和组件消融回答系统为什么逐步变好。第二层是外部机制对比，用 literature-inspired baselines 检验“单智能体一次性规划”或“仅靠执行反馈”是否已经足够。第三层是外部有效性验证，通过相邻技术子领域稳定性、软件工程数据上的跨领域小规模验证和人工盲评评估方法的跨数据集稳定性与人工可接受性。

\subsection{数据集与统一实验接口}

主数据集为 \texttt{main\_data\_security}，包含数据安全相关专利及其技术、法律、时间和组织字段。相邻技术子领域稳定性验证使用独立物化的数据集 \texttt{transfer\_iot\_auto}，用于检验方法在“数据安全专利主集 $\rightarrow$ 物联网与智能汽车安全子集”上的稳健性。真跨领域小验证使用 \texttt{cross\_domain\_bugzilla}，其字段至少覆盖 \texttt{title}、\texttt{description}、\texttt{severity}、\texttt{component}、\texttt{product/project}、\texttt{created\_at}、\texttt{resolved\_at}、\texttt{reporter} 和 \texttt{assignee}。三类数据集均通过统一注册表进行加载和调度，避免学位论文证据与期刊稿证据相互覆盖。

所有实验共享统一输入接口：\texttt{question}、\texttt{dataset\_profile}、\texttt{graph\_profile}、\texttt{knowledge\_enabled}、\texttt{iteration\_rounds} 和 \texttt{budget\_config}。输出接口固定为结构化计划 JSON、执行产物、自动化指标、语义评分和统计汇总。统一接口的作用是保证不同模式、不同数据集和不同证据池中的结果可以在同一评价协议下比较。

\subsection{分层实验设置}

\paragraph{主实验：渐进叠加}
主实验比较 \texttt{D\_baseline0}、\texttt{A\_template}、\texttt{B\_data\_aware} 和 \texttt{C\_iterative}。其中，\texttt{D\_baseline0} 代表纯 LLM 规划，\texttt{A\_template} 在其上增加知识模板，\texttt{B\_data\_aware} 进一步注入数据预览与证据锚点，\texttt{C\_iterative} 则在此基础上加入两轮规划-执行-反馈闭环。增强版正式目标规模为 \texttt{4 模式 × 5 问题 × 5 重复}，该组实验回答“系统为何逐步变好”。

\paragraph{组件消融}
组件消融比较 \texttt{C\_iterative}、\texttt{E\_ablate\_data} 和 \texttt{F\_ablate\_kg}，分别用于检验移除数据感知和移除知识图谱后的性能变化。该组实验回答“完整系统中核心组件各自贡献了什么”。主实验与消融共享同一主证据池，增强版正式目标为 \texttt{6 模式 × 5 问题 × 5 重复 = 150} 组内部证据；当前已冻结的可用内部证据仍为 \texttt{6 模式 × 3 问题 × 5 重复 = 90} 组。

\paragraph{外部基线}
外部基线采用两条文献启发的机制路线：\texttt{G\_react\_single\_agent}（ReAct-style single-agent baseline）和 \texttt{H\_execution\_feedback}（execution-feedback baseline）。两类基线均定位为 external mechanism baselines，而非严格文献系统复现。增强版正式目标是在主数据集上按 \texttt{5 问题 × 5 重复} 完成 \texttt{50} 组新增实验；当前已完成的版本为 \texttt{3 问题 × 5 重复 = 30} 组。

为避免事后倒选，子领域稳定性验证所使用的 strongest external mechanism baseline 在主数据集外部机制基线完成后即固定锁定，选择规则依次为：LLM 总分均值更高、LLM 总分标准差更小、\texttt{specificity\_references} 均值更高。当前已锁定的模式为 \texttt{H\_execution\_feedback}。

\paragraph{子领域稳定性验证}
相邻技术子领域稳定性验证固定在 \texttt{transfer\_iot\_auto} 上比较 \texttt{D\_baseline0}、\texttt{B\_data\_aware}、\texttt{C\_iterative} 和 strongest external mechanism baseline，正式目标为 \texttt{5 问题 × 3 重复 = 60} 组。这里的定位是“子领域稳定性检验”，其作用是说明方法不易过拟合主数据集；本文明确不把这一设置写成强跨领域迁移验证。此前的 \texttt{36} 组旧批次由于大量失败已整体废弃，不进入正文统计。

\paragraph{跨领域初步验证}
真跨领域验证使用 \texttt{cross\_domain\_bugzilla}，仅比较 \texttt{D\_baseline0} 与 \texttt{C\_iterative} 两个模式，按软件工程适配后的三类问题各重复 3 次，总规模为 \texttt{18} 组。该实验不复用专利知识图谱；若知识资源无法直接迁移，则以 \texttt{knowledge\_enabled=False} 运行，并将结果解释为框架的降级鲁棒性，而不是知识图谱的直接迁移能力。

\subsection{评价协议与统计方法}

本文采用双轨评价。自动化指标用于衡量计划是否能够覆盖关键字段、组织足够的方法组合并保持较高执行成功率；LLM-as-Judge 则从相关性、针对性、方法合理性、创新性和深度五个维度对匿名化方案进行评分。正文和表图统一报告重复实验的均值与标准差，并将 95\% bootstrap 置信区间作为关键配对比较的补充不确定性估计。

统计检验采用非参数路线。整体比较使用 Friedman 检验；关键配对比较使用 Wilcoxon signed-rank test 并进行 Holm 校正；效应量使用 rank-biserial correlation，并为关键比较补充 95\% bootstrap 置信区间。为保证口径一致，主实验、外部基线、子领域稳定性和跨领域验证分别在各自证据池内汇总，再按预定义模式集合组织统计结果，禁止将不同证据池中的原始结果直接混成单一总表。

\subsection{人工盲评协议}

人工盲评对象固定为匿名化后的规划方案 blueprint，不包含执行代码。盲评集合覆盖三类比较：\texttt{D/B/C}、\texttt{C/E/F} 和 \texttt{C/strongest external mechanism baseline}。每位评审对每份匿名方案按相关性、针对性、方法合理性、创新性、深度和可执行性六个维度给出 1--5 分。增强版投稿路线默认邀请 4 位评审，确保至少 3 位完成有效回收，并报告评审间一致性指标。匿名标签与模式映射分离保存，以降低先验偏见对评分的影响。

\subsection{证据目录与可追溯性}

为避免学位论文目录被覆盖，本文将主实验、外部基线、子领域稳定性验证、跨领域验证和人工盲评分别保存到独立证据目录中。正文表图只从对应目录导出的汇总文件中取数：内部机制归因使用 \texttt{\detokenize{outputs/repeated_experiments_full/}}，外部机制基线使用 \texttt{\detokenize{outputs/repeated_experiments_external_main/}}，子领域稳定性使用 \texttt{\detokenize{outputs/repeated_experiments_transfer_iot_auto/}}，跨领域验证使用 \texttt{\detokenize{outputs/repeated_experiments_cross_domain_bugzilla/}}，人工盲评使用 \texttt{\detokenize{outputs/jos_artifacts/human_review/}}。这一设计保证每个结论都可以追溯到独立、稳定且不互相覆盖的证据池。

\section{实验结果与讨论}

\subsection{六模式总览}

\begin{figure}[H]
\centering
\safeincludegraphics[0.92\textwidth]{fig4_six_mode_overview.png}
\caption{六模式总览（仅作全局观察）}
\label{fig:six-overview}
\end{figure}

六模式总览显示，\texttt{C\_iterative} 在当前重复实验中保持最高平均总分（23.87 $\pm$ 1.77），\texttt{E\_ablate\_data} 与 \texttt{F\_ablate\_kg} 处于第二梯队，\texttt{B\_data\_aware} 位于中间位置，而 \texttt{A\_template} 与 \texttt{D\_baseline0} 相对较低且彼此接近。需要强调的是，本图只用于全局观察，不承担机制归因；真正的归因仍应依赖四方案主实验和组件消融。

\subsection{主实验：四方案渐进叠加}

\begin{figure}[H]
\centering
\safeincludegraphics[0.82\textwidth]{fig3_main_total_scores.png}
\caption{四方案主实验总分对比}
\label{fig:main-total}
\end{figure}

\begin{figure}[H]
\centering
\safeincludegraphics[0.72\textwidth]{fig5_main_radar.png}
\caption{四方案五维评分雷达图}
\label{fig:main-radar}
\end{figure}

\IfFileExists{generated/table_main_llm.tex}{\begin{table}[htbp]
\centering
\caption{四方案 LLM-as-Judge 评分对比（5 问题 $\times$ 5 次重复的均值 $\pm$ 标准差）}
\label{tab:jos-main-llm}
\resizebox{\textwidth}{!}{%
\begin{tabular}{lllll}
\toprule
维度 & D\_baseline0 & A\_template & B\_data\_aware & C\_iterative \\
\midrule
相关性 & 3.68 ± 0.48 & 3.72 ± 0.54 & 4.04 ± 0.84 & 4.80 ± 0.41 \\
针对性 & 3.28 ± 0.46 & 3.64 ± 0.64 & 3.96 ± 0.73 & 4.80 ± 0.50 \\
方法合理性 & 3.44 ± 0.58 & 3.40 ± 0.65 & 3.64 ± 0.81 & 4.76 ± 0.52 \\
创新性 & 2.76 ± 0.60 & 2.92 ± 0.81 & 3.76 ± 0.60 & 4.68 ± 0.56 \\
深度 & 3.00 ± 0.50 & 2.88 ± 0.44 & 3.40 ± 0.76 & 4.76 ± 0.52 \\
总分(/25) & 16.16 ± 1.86 & 16.56 ± 2.36 & 18.80 ± 3.03 & 23.80 ± 2.12 \\
\bottomrule
\end{tabular}%
}
\end{table}
}{\par 主实验 LLM 汇总表待生成。}
\IfFileExists{generated/table_main_auto.tex}{\begin{table}[htbp]
\centering
\caption{四方案关键自动化指标对比（5 问题 $\times$ 5 次重复的均值 $\pm$ 标准差）}
\label{tab:jos-main-auto}
\resizebox{\textwidth}{!}{%
\begin{tabular}{lllll}
\toprule
指标 & D\_baseline0 & A\_template & B\_data\_aware & C\_iterative \\
\midrule
数据列覆盖率 & 0.264 ± 0.066 & 0.251 ± 0.060 & 0.275 ± 0.068 & 0.375 ± 0.084 \\
方法多样性 & 2.720 ± 0.542 & 2.800 ± 0.577 & 2.720 ± 0.458 & 5.120 ± 0.666 \\
针对性引用数 & 0.920 ± 1.115 & 1.880 ± 2.333 & 8.440 ± 4.744 & 14.080 ± 5.492 \\
参数丰富度 & 5.760 ± 1.899 & 5.000 ± 1.732 & 4.960 ± 2.091 & 11.440 ± 2.815 \\
生成任务总数 & 2.880 ± 0.526 & 2.880 ± 0.526 & 2.800 ± 0.408 & 5.160 ± 0.746 \\
代码执行成功率 & 1.000 ± 0.000 & 1.000 ± 0.000 & 1.000 ± 0.000 & 1.000 ± 0.000 \\
\bottomrule
\end{tabular}%
}
\vspace{0.3em}
\begin{minipage}{0.96\linewidth}
\footnotesize\raggedright
注：本表只保留与规划质量最直接相关的核心结构指标；完整指标集合见补充材料。
\end{minipage}
\end{table}
}{\par 主实验自动化指标表待生成。}
\IfFileExists{generated/table_question_main.tex}{\begin{table}[htbp]
\centering
\caption{四方案问题级 LLM 总分明细（均值 $\pm$ 标准差）}
\label{tab:jos-question-main}
\resizebox{\textwidth}{!}{%
\begin{tabular}{lllll}
\toprule
研究问题 & D\_baseline0 & A\_template & B\_data\_aware & C\_iterative \\
\midrule
Q1 技术影响力驱动因素 & 17.20 ± 2.39 & 17.40 ± 1.14 & 19.80 ± 2.39 & 23.20 ± 2.49 \\
Q2 技术融合趋势识别 & 16.40 ± 1.34 & 17.20 ± 3.42 & 18.00 ± 3.81 & 24.80 ± 0.45 \\
Q3 国家竞争态势评估 & 14.60 ± 1.34 & 16.40 ± 2.19 & 18.20 ± 1.30 & 23.40 ± 3.58 \\
Q4 技术生命周期阶段判断 & 15.60 ± 1.82 & 17.00 ± 1.87 & 16.40 ± 3.36 & 23.20 ± 2.17 \\
Q5 专利组合风险评估 & 17.00 ± 1.58 & 14.80 ± 2.59 & 21.60 ± 1.52 & 24.40 ± 0.55 \\
\bottomrule
\end{tabular}%
}
\end{table}
}{\par 主实验问题级表待生成。}

主实验反映出两个稳定结论。第一，\texttt{B\_data\_aware} 相比 \texttt{A\_template} 和 \texttt{D\_baseline0} 能够稳定提升针对性和总分，说明数据感知使系统从“问题语义驱动”转向“数据证据驱动”。第二，\texttt{C\_iterative} 相比 \texttt{B\_data\_aware} 在五个语义维度和多个结构指标上都进一步拉开差距，表明两轮迭代并不是简单堆叠任务数量，而是实质提高了方案的深度和方法组织能力。

值得注意的是，当前重复实验并不支持将 \texttt{A\_template} 稳定写成优于 \texttt{D\_baseline0} 的强结论。\texttt{A\_template} 确实在结构性上提供了一定先验，但如果缺乏数据感知和执行反馈，其收益仍不足以形成稳定梯度。

\subsection{组件消融}

\begin{figure}[H]
\centering
\safeincludegraphics[0.88\textwidth]{fig6_ablation_total.png}
\caption{组件消融结果：总分与针对性引用数}
\label{fig:ablation-total}
\end{figure}

\IfFileExists{generated/table_ablation_llm.tex}{\begin{table}[htbp]
\centering
\caption{组件消融 LLM-as-Judge 评分对比（5 问题 $\times$ 5 次重复的均值 $\pm$ 标准差）}
\label{tab:jos-ablation-llm}
\resizebox{\textwidth}{!}{%
\begin{tabular}{llll}
\toprule
维度 & C\_iterative & E\_ablate\_data & F\_ablate\_kg \\
\midrule
相关性 & 4.80 ± 0.41 & 4.72 ± 0.46 & 4.80 ± 0.41 \\
针对性 & 4.80 ± 0.50 & 4.32 ± 0.56 & 4.84 ± 0.47 \\
方法合理性 & 4.76 ± 0.52 & 4.44 ± 0.51 & 4.44 ± 0.71 \\
创新性 & 4.68 ± 0.56 & 3.72 ± 0.46 & 4.72 ± 0.54 \\
深度 & 4.76 ± 0.52 & 4.52 ± 0.51 & 4.32 ± 0.63 \\
总分(/25) & 23.80 ± 2.12 & 21.72 ± 2.01 & 23.12 ± 2.17 \\
\bottomrule
\end{tabular}%
}
\end{table}
}{\par 消融实验 LLM 汇总表待生成。}
\IfFileExists{generated/table_ablation_auto.tex}{\begin{table}[htbp]
\centering
\caption{组件消融关键自动化指标对比（5 问题 $\times$ 5 次重复的均值 $\pm$ 标准差）}
\label{tab:jos-ablation-auto}
\resizebox{\textwidth}{!}{%
\begin{tabular}{llll}
\toprule
指标 & C\_iterative & E\_ablate\_data & F\_ablate\_kg \\
\midrule
数据列覆盖率 & 0.375 ± 0.084 & 0.342 ± 0.085 & 0.365 ± 0.060 \\
方法多样性 & 5.120 ± 0.666 & 4.440 ± 0.583 & 5.080 ± 0.702 \\
针对性引用数 & 14.080 ± 5.492 & 3.800 ± 3.253 & 14.960 ± 4.641 \\
参数丰富度 & 11.440 ± 2.815 & 10.560 ± 2.063 & 12.280 ± 2.072 \\
生成任务总数 & 5.160 ± 0.746 & 4.680 ± 0.557 & 5.080 ± 0.400 \\
代码执行成功率 & 1.000 ± 0.000 & 1.000 ± 0.000 & 1.000 ± 0.000 \\
\bottomrule
\end{tabular}%
}
\end{table}
}{\par 消融实验自动化指标表待生成。}
\IfFileExists{generated/table_question_ablation.tex}{\begin{table}[htbp]
\centering
\caption{组件消融问题级 LLM 总分明细（均值 $\pm$ 标准差）}
\label{tab:jos-question-ablation}
\resizebox{\textwidth}{!}{%
\begin{tabular}{llll}
\toprule
研究问题 & C\_iterative & E\_ablate\_data & F\_ablate\_kg \\
\midrule
Q1 技术影响力驱动因素 & 23.20 ± 2.49 & 21.40 ± 2.30 & 24.80 ± 0.45 \\
Q2 技术融合趋势识别 & 24.80 ± 0.45 & 23.40 ± 0.55 & 22.80 ± 1.48 \\
Q3 国家竞争态势评估 & 23.40 ± 3.58 & 22.20 ± 1.79 & 23.00 ± 1.87 \\
Q4 技术生命周期阶段判断 & 23.20 ± 2.17 & 20.20 ± 1.79 & 23.40 ± 2.30 \\
Q5 专利组合风险评估 & 24.40 ± 0.55 & 21.40 ± 2.30 & 21.60 ± 3.21 \\
\bottomrule
\end{tabular}%
}
\end{table}
}{\par 消融实验问题级表待生成。}

组件消融进一步说明了完整系统的性能来源。对比 \texttt{C\_iterative} 与 \texttt{E\_ablate\_data} 可见，移除数据感知后，针对性引用数由 14.20 降至 3.00，总分也同步下降。这说明数据感知并不是简单增加上下文长度，而是直接改变规划起点。对比 \texttt{C\_iterative} 与 \texttt{F\_ablate\_kg} 则可见，知识图谱在平均意义上仍提供正向收益，但该收益具有任务依赖性：在开放度更高的问题上，较弱的图谱约束可能释放探索空间；而在结构锚点清晰、需要链式推理的问题上，图谱对变量组织与方法选择的约束更能带来稳定优势。

\subsection{外部机制基线比较}

\IfFileExists{figures/fig7_external_total.png}{
\begin{figure}[H]
\centering
\safeincludegraphics[0.84\textwidth]{fig7_external_total.png}
\caption{完整系统与外部机制基线总分对比}
\label{fig:external-total}
\end{figure}
}{}

\IfFileExists{generated/table_external_llm.tex}{\begin{table}[htbp]
\centering
\caption{完整系统与外部基线的 LLM-as-Judge 评分对比（均值 $\pm$ 标准差）}
\label{tab:jos-external-llm}
\resizebox{\textwidth}{!}{%
\begin{tabular}{llll}
\toprule
维度 & C\_iterative & G\_react\_single\_agent & H\_execution\_feedback \\
\midrule
相关性 & 4.80 ± 0.41 & 3.96 ± 0.61 & 4.68 ± 0.80 \\
针对性 & 4.80 ± 0.50 & 3.40 ± 0.82 & 4.56 ± 0.87 \\
方法合理性 & 4.76 ± 0.52 & 3.28 ± 0.68 & 4.16 ± 0.80 \\
创新性 & 4.68 ± 0.56 & 2.16 ± 0.55 & 3.76 ± 0.52 \\
深度 & 4.76 ± 0.52 & 2.36 ± 0.81 & 4.08 ± 0.81 \\
总分(/25) & 23.80 ± 2.12 & 15.16 ± 3.09 & 21.24 ± 3.26 \\
\bottomrule
\end{tabular}%
}
\end{table}
}{\par 外部基线 LLM 汇总表待生成。}
\IfFileExists{generated/table_external_auto.tex}{\begin{table}[htbp]
\centering
\caption{完整系统与外部基线的关键自动化指标对比（均值 $\pm$ 标准差）}
\label{tab:jos-external-auto}
\resizebox{\textwidth}{!}{%
\begin{tabular}{llll}
\toprule
指标 & C\_iterative & G\_react\_single\_agent & H\_execution\_feedback \\
\midrule
数据列覆盖率 & 0.375 ± 0.084 & 0.200 ± 0.075 & 0.240 ± 0.076 \\
方法多样性 & 5.120 ± 0.666 & 1.120 ± 0.440 & 1.280 ± 0.678 \\
针对性引用数 & 14.080 ± 5.492 & 1.320 ± 1.994 & 2.120 ± 2.351 \\
参数丰富度 & 11.440 ± 2.815 & 2.040 ± 2.226 & 5.600 ± 3.571 \\
生成任务总数 & 5.160 ± 0.746 & 2.560 ± 0.651 & 4.560 ± 0.583 \\
代码执行成功率 & 1.000 ± 0.000 & 1.000 ± 0.000 & 1.000 ± 0.000 \\
\bottomrule
\end{tabular}%
}
\end{table}
}{\par 外部基线自动化指标表待生成。}
\IfFileExists{generated/table_external_question.tex}{\begin{table}[htbp]
\centering
\caption{完整系统与外部基线的问题级 LLM 总分明细（均值 $\pm$ 标准差）}
\label{tab:jos-external-question}
\resizebox{\textwidth}{!}{%
\begin{tabular}{llll}
\toprule
研究问题 & C\_iterative & G\_react\_single\_agent & H\_execution\_feedback \\
\midrule
Q1 技术影响力驱动因素 & 23.20 ± 2.49 & 13.80 ± 0.45 & 23.20 ± 0.84 \\
Q2 技术融合趋势识别 & 24.80 ± 0.45 & 16.00 ± 3.81 & 18.20 ± 5.21 \\
Q3 国家竞争态势评估 & 23.40 ± 3.58 & 14.80 ± 4.09 & 21.00 ± 3.00 \\
Q4 技术生命周期阶段判断 & 23.20 ± 2.17 & 14.80 ± 1.30 & 22.40 ± 1.34 \\
Q5 专利组合风险评估 & 24.40 ± 0.55 & 16.40 ± 4.34 & 21.40 ± 2.61 \\
\bottomrule
\end{tabular}%
}
\end{table}
}{\par 外部基线问题级表待生成。}
\IfFileExists{generated/table_external_significance.tex}{\begin{table}[htbp]
\centering
\caption{完整系统与外部机制基线的 Wilcoxon 检验结果（LLM 总分）}
\label{tab:jos-external-significance}
\resizebox{\textwidth}{!}{%
\begin{tabular}{lllllll}
\toprule
比较 & n & 均值差 & 95\% CI & p & Holm校正后p & 效应量 \\
\midrule
C\_iterative vs G\_react\_single\_agent & 25 & 8.64 & [7.32, 9.92] & 0.0000 & 0.0000 & 1.000 \\
C\_iterative vs H\_execution\_feedback & 25 & 2.56 & [1.00, 4.24] & 0.0037 & 0.0037 & 0.688 \\
\bottomrule
\end{tabular}%
}
\vspace{0.3em}
\begin{minipage}{0.96\linewidth}
\footnotesize\raggedright
注：均值差按“左侧模式减右侧模式”计算，效应量为 rank-biserial correlation。
\end{minipage}
\end{table}
}{\par 外部基线显著性表待生成。}

\IfFileExists{generated/table_external_llm.tex}{
外部机制基线节用于回答一个不同于内部消融的问题：如果只采用单智能体一次性规划，或只采用执行反馈驱动再规划，而不显式引入数据感知与知识约束闭环，系统是否已经能够达到与完整方法相近的稳定性。当前已完成的 30 组外部机制基线已锁定 \texttt{H\_execution\_feedback} 为当前 strongest external mechanism baseline，且 \texttt{C\_iterative} 在总分均值和稳定性上均优于 \texttt{G/H}。正文在这一节只引用外部证据池及与之对齐的 \texttt{C\_iterative} 汇总，不把主实验和外部机制基线的原始结果直接混成一个未分层的大表。正式数值以图~\ref{fig:external-total} 及外部机制基线表为准；5 问题口径的最终结果将在增强版实验完成后替换当前 3 问题版本。
}{
本节已经预留外部机制基线图表与显著性表接口。若当前编译版本尚未显示对应表图，表示 \texttt{\detokenize{outputs/repeated_experiments_external_main/}} 的正式汇总或导出仍在进行中；在结果落地之前，正文不对 \texttt{C\_iterative} 与 \texttt{G/H} 的优劣关系做超出数据的书面结论。
}

\subsection{子领域稳定性验证}

\IfFileExists{figures/fig8_transfer_total.png}{
\begin{figure}[H]
\centering
\safeincludegraphics[0.84\textwidth]{fig8_transfer_total.png}
\caption{子领域稳定性验证总分对比}
\label{fig:transfer-total}
\end{figure}
}{}

\IfFileExists{generated/table_transfer_llm.tex}{\input{generated/table_transfer_llm.tex}}{\par 子领域稳定性验证 LLM 汇总表待生成。}
\IfFileExists{generated/table_transfer_auto.tex}{\input{generated/table_transfer_auto.tex}}{\par 子领域稳定性验证自动化指标表待生成。}
\IfFileExists{generated/table_transfer_question.tex}{\input{generated/table_transfer_question.tex}}{\par 子领域稳定性验证问题级表待生成。}
\IfFileExists{generated/table_transfer_significance.tex}{\input{generated/table_transfer_significance.tex}}{\par 子领域稳定性验证显著性表待生成。}

\IfFileExists{generated/table_transfer_llm.tex}{
子领域稳定性验证用于检验完整系统的优势是否能在相邻技术子领域中保持。该节只比较 \texttt{D\_baseline0}、\texttt{B\_data\_aware}、\texttt{C\_iterative} 和 strongest external mechanism baseline，不再引入全部六模式，从而把讨论聚焦到“完整系统是否仍是最稳选择”这一问题。若子领域数据集中的模式排序与主数据集一致，则说明本文方法并非只对单一主数据集过拟合；若出现局部偏移，则需要结合问题级结果分析不同子领域中的证据结构差异。需要强调的是，这一设置只能支撑“子领域稳定性”，不支撑强跨领域泛化。
}{
当前旧批次 \texttt{transfer\_iot\_auto} 结果因大量失败已整体作废，不能进入正文。增强版路线将以 \texttt{D/B/C/H}、5 问题、3 重复的 60 组设计重跑该证据池，并仅将其写作“子领域稳定性验证”；在正式结果落地前，正文不提前写入优劣结论。
}

\subsection{跨领域初步验证}

\IfFileExists{figures/fig9_cross_domain_total.png}{
\begin{figure}[H]
\centering
\safeincludegraphics[0.78\textwidth]{fig9_cross_domain_total.png}
\caption{跨领域初步验证总分对比}
\label{fig:cross-domain-total}
\end{figure}
}{}

\IfFileExists{generated/table_cross_domain_llm.tex}{\begin{table}[htbp]
\centering
\caption{跨领域初步验证的 LLM-as-Judge 评分对比（均值 $\pm$ 标准差）}
\label{tab:jos-cross-domain-llm}
\resizebox{\textwidth}{!}{%
\begin{tabular}{lll}
\toprule
维度 & D\_baseline0 & C\_iterative \\
\midrule
相关性 & 4.00 ± 0.50 & 4.78 ± 0.67 \\
针对性 & 3.22 ± 0.44 & 5.00 ± 0.00 \\
方法合理性 & 3.33 ± 0.87 & 4.89 ± 0.33 \\
创新性 & 2.22 ± 0.67 & 4.89 ± 0.33 \\
深度 & 3.11 ± 0.78 & 4.89 ± 0.33 \\
总分(/25) & 15.89 ± 2.93 & 24.44 ± 1.33 \\
\bottomrule
\end{tabular}%
}
\end{table}
}{\par 跨领域初步验证 LLM 汇总表待生成。}
\IfFileExists{generated/table_cross_domain_question.tex}{\begin{table}[htbp]
\centering
\caption{跨领域初步验证的问题级 LLM 总分明细（均值 $\pm$ 标准差）}
\label{tab:jos-cross-domain-question}
\resizebox{\textwidth}{!}{%
\begin{tabular}{lll}
\toprule
研究问题 & D\_baseline0 & C\_iterative \\
\midrule
Q1 缺陷影响因素分析 & 18.00 ± 4.36 & 23.67 ± 2.31 \\
Q2 缺陷与模块演化趋势 & 14.33 ± 2.08 & 25.00 ± 0.00 \\
Q3 项目/组件质量评估 & 15.33 ± 0.58 & 24.67 ± 0.58 \\
\bottomrule
\end{tabular}%
}
\end{table}
}{\par 跨领域初步验证问题级表待生成。}

\IfFileExists{generated/table_cross_domain_llm.tex}{
跨领域初步验证用于回答一个更严格的问题：当数据从专利分析切换到软件工程 Bugzilla 缺陷数据时，完整系统相较纯 LLM 基线是否仍然保持优势。由于该部分只比较 \texttt{D\_baseline0} 与 \texttt{C\_iterative}，且默认不复用专利知识图谱，因此它更适合作为“框架降级后是否仍稳健”的证据，而不是对知识图谱跨领域迁移能力的直接证明。
}{
跨领域初步验证的注册表、数据物化脚本和表图接口已经接入工作区，但当前尚无可用 Bugzilla 证据。正式稿中，这一节只会在真实 18 组实验完成后写入结果，不会用主数据集或子领域数据代替跨领域证据。
}

\subsection{统计结果}

\IfFileExists{generated/table_significance.tex}{\begin{table}[htbp]
\centering
\caption{关键配对比较的 Wilcoxon 检验结果（LLM 总分）}
\label{tab:jos-significance}
\resizebox{\textwidth}{!}{%
\begin{tabular}{lllllll}
\toprule
比较 & n & 均值差 & 95\% CI & p & Holm校正后p & 效应量 \\
\midrule
B\_data\_aware vs A\_template & 15 & 2.07 & [0.00, 4.07] & 0.0780 & 0.3119 & 0.533 \\
B\_data\_aware vs D\_baseline0 & 15 & 3.13 & [1.40, 4.60] & 0.0084 & 0.0418 & 0.758 \\
C\_iterative vs B\_data\_aware & 15 & 5.27 & [3.53, 7.00] & 0.0003 & 0.0027 & 0.958 \\
C\_iterative vs E\_ablate\_data & 15 & 1.80 & [0.87, 2.80] & 0.0055 & 0.0331 & 0.939 \\
C\_iterative vs F\_ablate\_kg & 15 & 1.67 & [-0.40, 3.80] & 0.1754 & 0.3508 & 0.410 \\
\bottomrule
\end{tabular}%
}
\vspace{0.3em}
\begin{minipage}{0.96\linewidth}
\footnotesize\raggedright
注：均值差按“左侧模式减右侧模式”计算，效应量为 rank-biserial correlation。
\end{minipage}
\end{table}
}{\par 主实验显著性表待生成。}

当前内部证据已经补充 Friedman 与 Wilcoxon 检验，并为关键配对比较报告 bootstrap 置信区间和 rank-biserial correlation。统计检验的作用不是把论文写成“所有差异都显著”的故事，而是把当前可发表结论限定在真正有重复支持的范围内。就已冻结的 90 组主证据而言，最稳的结论依然是：数据感知提供稳定增益，两轮迭代决定性能上限，知识图谱平均上有帮助但收益依赖任务结构。

\subsection{人工盲评与讨论}

\IfFileExists{figures/fig10_human_review.png}{
\begin{figure}[H]
\centering
\safeincludegraphics[0.82\textwidth]{fig10_human_review.png}
\caption{人工盲评总体评分对比}
\label{fig:human-review}
\end{figure}
}{}

\IfFileExists{generated/table_human_review.tex}{\input{generated/table_human_review.tex}}{\par 人工盲评汇总表待生成。}

\IfFileExists{generated/table_human_review.tex}{
人工盲评用于补足 LLM-as-Judge 的外部视角。由于盲评对象被限制为匿名化 blueprint 而非执行代码，这一结果更接近人类评审者对“规划质量本身”的判断，而不是对脚本调试能力的判断。若人工盲评与自动汇总方向一致，则说明完整系统的优势不仅体现在模型互评分数上，也体现在人类对方案针对性、深度和可执行性的直观认可上。
}{
人工盲评材料包、匿名映射和评审模板已经接入工作区。若当前编译版本尚未显示盲评分数，表示真实评审表尚未回收或尚未汇总；这一状态会在有效性威胁中明确说明，而不会被写成已经完成的人类验证。
}

从机制上看，三项能力分别对应不同层级的问题。数据感知回答“当前数据最值得看什么”，知识图谱回答“面对这类证据通常该怎么分析”，两轮迭代回答“首轮方案不足时如何继续深化”。如果只保留其中任意两个组件，系统仍可能在某些问题上取得较高分数；但若要在多类问题上同时保持高针对性、高方法合理性与高深度，完整配置仍是目前证据支持下最稳妥的选择。

\section{有效性威胁}

本文的主要有效性威胁体现在以下五个方面。

\begin{itemize}
  \item \textbf{外部机制基线的实现仍属文献启发，而非严格复现。} 本文引入的 ReAct-style single-agent baseline 和 execution-feedback baseline 用于覆盖代表性机制路线，但它们并不声称逐字复现某篇论文的完整系统。因此，外部比较更适合支撑“机制层面对比”，而不是“逐篇文献胜负”。
  \item \textbf{子领域稳定性不等于强跨领域泛化。} \texttt{transfer\_iot\_auto} 与主数据集同属专利数据，只能说明系统在相邻技术子领域上的稳定性，不能直接支撑“跨领域迁移能力”。真正的跨领域主张必须建立在 Bugzilla 等异质数据上的独立验证之上。
  \item \textbf{数据集与问题覆盖仍有限。} 即便增强版路线把主问题扩展到五类，并补充软件工程跨领域小验证，任务集合与跨领域样本规模仍然有限。未来若面向更宽泛的软件工程或科学分析场景推广，仍需在更多数据结构、更多任务类型和更强异质性数据集上继续验证。
  \item \textbf{语义评价仍存在模型偏好与人工样本规模限制。} LLM-as-Judge 能够提供稳定的批量比较，但其评分标准仍可能受模型偏好影响\citep{zheng2023judging,ji2023survey}；人工盲评则更接近人类评审视角，但样本规模、评审背景和时间成本会限制其统计力度。因此，两类评价只能相互补充，不能互相替代。
  \item \textbf{知识约束强度尚未实现自适应调节。} 当前实验已经表明知识图谱收益具有任务依赖性，但本文尚未实现根据问题开放度、证据密度和方法可行性自动调节图谱约束强度的策略。这意味着当前框架虽然能够识别“何时可能过约束”，但还不能自动完成更细粒度的强弱切换。
\end{itemize}

\IfFileExists{generated/table_external_llm.tex}{
外部机制基线结果已经被纳入正文后，第一类威胁由“缺少外部对比”收缩为“外部对比并非严格文献复现”；若当前编译版本尚未显示外部表图，则该缺失本身仍构成投稿前必须显式说明的证据限制。
}{
若当前编译版本尚未显示外部机制基线、子领域稳定性、跨领域验证或人工盲评表图，则这些证据池仍处于回填阶段；这种缺失会直接削弱“外部对比”和“外部有效性”的说服力，因此不能被模糊表述为已经完成。
}

这些威胁不会推翻当前已经观察到的机制趋势，但它们要求本文在写作上始终保持口径克制：内部主张只由已冻结的主证据支持，外部结论只在外部证据真正落地后写入正文，人类评价只在真实盲评分数回收后进入最终结论。

\section{结论}

本文提出了一种面向结构化分析任务的数据感知多智能体规划框架，并在专利分析场景中构建了独立于学位论文的期刊稿证据工作区。与只优化执行层的系统不同，本文把关注点前移到规划层，并通过 Evidence Grounding、Knowledge-Constrained Planning 与 Iterative Refinement 三项机制，把数据证据、知识约束和执行反馈组织到统一闭环中。

基于已冻结的 90 组内部重复实验与已完成的 30 组外部机制基线，本文已经可以得到三项稳健判断：其一，数据感知是推动规划从问题语义驱动走向证据驱动的关键机制；其二，两轮迭代决定了系统可达到的性能上限；其三，知识图谱在平均意义上提供正向贡献，但收益强度受问题结构显著影响。换言之，完整系统当前在内部机制归因层面已经表现出“总体最优且相对稳定”的特征，而 \texttt{B\_data\_aware} 则稳定优于不含数据感知的方案。

\IfFileExists{generated/table_external_llm.tex}{
外部机制基线已经说明，仅靠单智能体 ReAct 风格规划或仅靠执行反馈，并不足以达到完整系统的稳定表现。接下来的增强版任务是把这一判断扩展到五问题口径、子领域稳定性、Bugzilla 跨领域小验证和人工盲评，从而形成“内部归因 + 外部机制对比 + 外部有效性”的完整闭环。若这些补充证据与当前趋势保持一致，则本文可以进一步支持这样一个结论：本文提出的是一种具有外部可接受性的规划方法，而不只是针对单一评分器优化的实验系统。
}{
外部机制基线、子领域稳定性、跨领域验证和人工盲评接口已经接入期刊稿工作区，但若当前编译版本尚未显示对应表图，则相关证据仍在回填过程中。对此，本文保持审慎口径：在外部证据真正落地之前，不把“完整系统在外部对比和外部有效性验证中均最优”写成既成事实。
}

总体而言，本文把结构化分析系统的核心问题从“如何把代码跑出来”前移到“如何先形成更合理的分析计划”，并给出了一条可复验、可扩展、可逐层验证的软件方法路线。后续工作将继续围绕五问题扩展、相邻子领域稳定性重跑、Bugzilla 跨领域验证和更大规模人工盲评完善证据链，同时探索知识约束强度的自适应调节机制。


\bibliographystyle{unsrtnat}
\bibliography{references}

\end{document}
